% !TeX root=main.tex
% دستور زیر باید در اولین فصل شما باشد. آن را حذف نکنید!
\pagenumbering{arabic}

\thispagestyle{empty}

\chapter{مقدمه}

\section{بیان مسئله}

افزایش وضوح تصویر تک‌تصویری یکی از مسائل بنیادی در پردازش تصویر و بینایی ماشین است که طی دو دهه اخیر به‌دلیل گسترش کاربردهای عملی آن، توجه گسترده‌ای را به خود جلب کرده است. در بسیاری از سامانه‌های تصویربرداری، از جمله دوربین‌های تلفن همراه، سامانه‌های نظارتی، تصاویر تاریخی اسکن‌شده، تصاویر پزشکی و داده‌های سنجش از دور، محدودیت‌های سخت‌افزاری، فاصله از سوژه، محدودیت پهنای باند انتقال داده یا شرایط محیطی موجب می‌شود تنها نسخه‌ای با وضوح پایین از تصویر در دسترس باشد. در چنین شرایطی، بازسازی تصویری با وضوح بالاتر که از نظر ساختاری و ادراکی به واقعیت نزدیک باشد، اهمیت فراوانی می‌یابد.

مسئله ابرتفکیک تک‌تصویری ذاتاً بدوضع\LTRfootnote{Ill-posed} است. علت این امر آن است که فرآیند تولید تصویر با وضوح پایین، اطلاعات فرکانسی بالا را از بین می‌برد و بنابراین بازگرداندن این اطلاعات بدون در اختیار داشتن قیود پیشین مناسب، امکان‌پذیر نیست. مدل عمومی تخریب تصویر را می‌توان به‌صورت زیر بیان کرد:

\begin{equation}
	\mathbf{y} = (\mathbf{x} \otimes \mathbf{k}) \downarrow_s + \mathbf{n}
\end{equation}

که در آن $\mathbf{x}$ تصویر با وضوح بالا، $\mathbf{k}$ کرنل تاری، $\otimes$ عملگر کانولوشن، $\downarrow_s$ عملگر نمونه‌برداری با ضریب مقیاس $s$ و $\mathbf{n}$ نویز افزوده است. در قالبی خطی‌تر، این رابطه را می‌توان به شکل زیر نیز نوشت:

\begin{equation}
	\mathbf{y} = \mathbf{D}\mathbf{x} + \mathbf{n}
\end{equation}

که در آن $\mathbf{D}$ عملگر ترکیبی تاری و کاهش نمونه‌برداری است. هدف ابرتفکیک، تخمین $\mathbf{x}$ از روی مشاهده‌ی $\mathbf{y}$ است. با توجه به اینکه عملگر $\mathbf{D}$ معمولاً رتبه کامل ندارد و بخشی از اطلاعات را حذف می‌کند، این مسئله فاقد جواب یکتا بوده و نیازمند استفاده از پیشین‌های ساختاری یا آماری است.

یکی از رویکردهای مهم در این حوزه، بهره‌گیری از نمایش تنک\LTRfootnote{Sparse Representation} است. در این چارچوب فرض می‌شود هر قطعه کوچک از تصویر را می‌توان به‌صورت ترکیب خطی تعداد اندکی اتم از یک دیکشنری مناسب بیان کرد. اگر $\mathbf{y}$ یک پچ با وضوح پایین و $\mathbf{D}_l$ دیکشنری متناظر باشد، مسئله یافتن ضرایب تنک به شکل زیر بیان می‌شود:

\begin{equation}
	\boldsymbol{\alpha}^\star =
	\arg\min_{\boldsymbol{\alpha}}
	\left\|
	\mathbf{y} - \mathbf{D}_l \boldsymbol{\alpha}
	\right\|_2^2
	+
	\lambda
	\left\|
	\boldsymbol{\alpha}
	\right\|_1
\end{equation}

فرض کلیدی در این رویکرد آن است که همین ضرایب $\boldsymbol{\alpha}^\star$ را می‌توان برای بازسازی پچ متناظر در فضای وضوح بالا نیز به کار برد:

\begin{equation}
	\mathbf{x} = \mathbf{D}_h \boldsymbol{\alpha}^\star
\end{equation}

که در آن $\mathbf{D}_h$ دیکشنری فضای وضوح بالا است. این ایده نشان داد که تنکی یک پیشین قدرتمند برای مدل‌سازی ساختارهای طبیعی تصویر است. با این حال، حل مسائل بهینه‌سازی تنک برای تمامی پچ‌های تصویر، هزینه محاسباتی قابل توجهی دارد و در کاربردهای بزرگ‌مقیاس به چالشی جدی تبدیل می‌شود.

با پیشرفت یادگیری عمیق، شبکه‌های کانولوشنی بسیار عمیق که بر روی پایگاه‌داده‌های بزرگ آموزش دیده‌اند، عملکرد ابرتفکیک را به شکل چشمگیری بهبود بخشیدند. این روش‌ها به‌صورت داده‌محور، نگاشت میان تصویر با وضوح پایین و بالا را می‌آموزند و در شرایطی که مدل تخریب مشخص و ایده‌آل باشد، به نتایج درخشانی دست می‌یابند. در اغلب این مطالعات، تصویر با وضوح پایین از طریق یک کرنل نمونه‌برداری مشخص و بدون حضور نویز یا مصنوعات دیگر تولید می‌شود.

با این حال، همان‌گونه که در پژوهش‌های اخیر نشان داده شده است، تصاویر دنیای واقعی به‌ندرت از چنین مدل تخریب ایده‌آلی تبعیت می‌کنند. در عمل، تصویر ممکن است تحت تأثیر تاری لنز، خطاهای فوکوس، نویز حسگر، فشرده‌سازی با اتلاف، تقویت بیش‌ازحد لبه‌ها یا اشتراک‌گذاری چندباره در شبکه‌های اجتماعی قرار گرفته باشد. این فرآیند تخریب را می‌توان به‌صورت یک ترکیب چندمرحله‌ای مدل کرد:

\begin{equation}
	\mathbf{y} =
	\mathcal{C}_2
	\Big(
	\mathcal{D}_2
	(
	\mathcal{C}_1
	(
	\mathcal{D}_1(\mathbf{x})
	)
	)
	\Big)
	+
	\mathbf{n}
\end{equation}

که در آن $\mathcal{D}_i$ عملیات تاری و کاهش نمونه‌برداری و $\mathcal{C}_i$ فرآیندهای فشرده‌سازی یا پردازش‌های غیرخطی را مدل می‌کنند. این دیدگاه که تحت عنوان تخریب مرتبه‌بالا\LTRfootnote{High-Order Degradation} مطرح شده است، نشان می‌دهد فاصله قابل توجهی میان داده‌های مصنوعی مورد استفاده در آموزش و تصاویر واقعی وجود دارد.

در پاسخ به این چالش، رویکردهای مبتنی بر شبکه‌های مولد تخاصمی\LTRfootnote{Generative Adversarial Networks} مطرح شدند که هدف آن‌ها نه تنها کاهش خطای بازسازی عددی، بلکه افزایش واقع‌گرایی ادراکی تصویر خروجی است. در این چارچوب، یک شبکه مولد تلاش می‌کند تصویر با وضوح بالا را بازسازی کند، در حالی که یک شبکه تمایزبخش کیفیت و طبیعی‌بودن خروجی را ارزیابی می‌کند. نسخه‌های توسعه‌یافته‌ای مانند روش تقویت‌شده برای تصاویر واقعی\LTRfootnote{Real-ESRGAN+} با در نظر گرفتن مدل‌های تخریب پیچیده‌تر و طراحی ساختارهای پایدارتر برای آموزش، سعی در کاهش فاصله میان داده‌های مصنوعی و تصاویر دنیای واقعی داشته‌اند. این رویکردها نشان می‌دهند که علاوه بر دقت بازسازی، کیفیت ادراکی و سازگاری با شرایط غیرایده‌آل نیز به مؤلفه‌ای کلیدی در ارزیابی سامانه‌های ابرتفکیک تبدیل شده است.

در سوی دیگر، رویکرد یادگیری درونی\LTRfootnote{Internal Learning} با تکیه بر خاصیت بازگشت‌پذیری چندمقیاسی در تصاویر طبیعی، پیشنهاد می‌کند که نمونه‌های آموزشی مورد نیاز برای یادگیری نگاشت وضوح پایین به بالا را می‌توان مستقیماً از خود تصویر استخراج کرد. این ایده مبنای روش ابرتفکیک بدون آموزش پیشین\LTRfootnote{Zero-Shot Super-Resolution} است که در آن یک شبکه کوچک در زمان آزمون و تنها با استفاده از نمونه‌های استخراج‌شده از تصویر ورودی آموزش داده می‌شود و به‌صورت تصویر-ویژه با شرایط تخریب سازگار می‌گردد.

در نتیجه، مسئله ابرتفکیک تک‌تصویری در ادبیات پژوهش را می‌توان در تعامل میان سه محور اساسی تحلیل کرد: بهره‌گیری از پیشین‌های ساختاری مانند تنکی، استفاده از مدل‌های یادگیری عمیق داده‌محور و توجه به واقع‌گرایی ادراکی و مدل‌سازی دقیق تخریب در تصاویر واقعی. با وجود پیشرفت‌های چشمگیر در هر یک از این محورها، همچنان چالش‌هایی نظیر پیچیدگی محاسباتی بالا، هزینه آموزش، و شکاف میان داده‌های مصنوعی و شرایط عملی وجود دارد. ازاین‌رو، بازنگری در ساختار محاسباتی و بهینه‌سازی روش‌های موجود، در حالی که بنیان نظری آن‌ها حفظ شود، گامی ضروری در جهت کاربردی‌سازی مؤثر این فناوری‌ها به شمار می‌آید.

\section{روش پیشنهادی}

در این پژوهش، تمرکز اصلی بر بهبود کارایی دو رویکرد مطرح در ابرتفکیک، یعنی روش مبتنی بر نمایش تنک و روش مبتنی بر یادگیری درونی بوده است.

در بهبود روش مبتنی بر نمایش تنک، چارچوب اصلی حفظ شده اما فرآیند حل مسئله بازنمایی تنک بازطراحی شده است. به جای استفاده از روش‌های بهینه‌سازی تکراری پرهزینه، از یک راهبرد انتخاب تدریجی اتم‌های مؤثر استفاده شده است که در هر مرحله مناسب‌ترین اتم را بر اساس شباهت به باقیمانده انتخاب می‌کند. همچنین ساختار استخراج ویژگی و نحوه پردازش قطعه‌های تصویر به‌گونه‌ای بازآرایی شده که محاسبات تکراری کاهش یابد و پیاده‌سازی کارآمدتری حاصل شود.

در بهبود روش مبتنی بر یادگیری درونی نیز چارچوب کلی آموزش شبکه تصویر-ویژه حفظ شده است، اما سازوکار کنترل فرآیند آموزش بازتنظیم شده است. با اصلاح سیاست توقف، بازطراحی نحوه نمونه‌برداری از تصویر و تنظیم راهبرد به‌روزرسانی نرخ یادگیری، فرآیند آموزش پایدارتر و کارآمدتر شده است، بدون آنکه اصول بنیادین یادگیری درونی تغییر یابد.

هدف از این اصلاحات، دستیابی به تعادلی مناسب میان کیفیت بازسازی و کارایی محاسباتی، در چارچوب همان مدل‌های پایه بوده است.

\section{ساختار پایان‌نامه}

این پایان‌نامه در پنج فصل تنظیم شده است. در فصل اول، بیان مسئله، انگیزه پژوهش و معرفی کلی روش پیشنهادی ارائه شد.  

فصل دوم به بررسی و تشریح مفاهیم پایه، مبانی نظری و پیشینه پژوهش اختصاص دارد. در این فصل، اصول نمایش تنک، مفاهیم یادگیری عمیق، ساختار شبکه‌های کانولوشنی، یادگیری درونی، مدل‌های تخریب تصویر و معیارهای ارزیابی کیفیت بازسازی به‌صورت نظام‌مند مرور می‌شوند تا چارچوب نظری لازم برای فصول بعد فراهم گردد.

در فصل سوم، روش‌های پایه مورد استفاده در این پژوهش به‌صورت دقیق معرفی و تحلیل می‌شوند و پیاده‌سازی آن‌ها مورد بررسی قرار می‌گیرد.  

فصل چهارم به ارائه بهبودهای پیشنهادی، تحلیل مقایسه‌ای و ارزیابی تجربی نتایج اختصاص دارد.  

در نهایت، فصل پنجم شامل جمع‌بندی، نتیجه‌گیری و پیشنهادهایی برای پژوهش‌های آینده است.